% Transcription — fonds Alexandre Grothendieck, shelfmark 120, batch 2 (pages 21-24).
% Established from the facsimile archives/batches/120/batch-02.pdf (a copy of
% 120.pdf fetched through the project's relay,
% https://grothendieck-relay.onrender.com/source/120.pdf; PDF page k is the
% archivists' page 20+k, no cover sheet), per the skill transcribe-grothendieck.
% The folder is twenty-four pages; this is the second and last batch (pages
% 1-20 are batch 1, transcribed separately and in parallel, so the notation
% below is fixed from these four pages alone).
%
% Pass: Opus 5.5 (claude-opus-5-5), 2026-09-26 — first pass, unchecked against the pages by a human.
%
% Pages skipped: none. Pages summarised: none. Pages transcribed: 21-24, all
% four in his hand, one ink, fast drafting; the formulas carry, several
% connective words do not and are left \ill{}.
%
% His pagination: page 22 carries « 11 » top right, presumably his own leaf
% number (recorded once in a \note); the other three leaves show none.
%
% What the batch is about. Tame topology (« topologie modérée », his term, as
% in Esquisse d'un programme), in an axiomatic set-up: a class \mathcal{M} of
% « parties modérées » of the R^n subject to numbered axioms (M1) ... (M11),
% of which (M6), (M7), (M8), (M11) are cited here, R_0 a subfield of R.
% Page 21: ind-tame spaces (R^I for an arbitrary index set I, when R is the
% union of its tame parts, (M8)), tame functions on them containing the
% polynomial rings R_0[(X_i)] in the coordinates; an interlinear addition on
% R_0-polynomial maps R^I -> X being tame. Corollary: a tame function on a
% tame space is invertible in the ring of tame functions iff it does not
% vanish; reduced to a second corollary, that x -> x^{-1} restricted to a
% tame A in R^* is tame. Page 22: the proof (the graph as the preimage of 1
% under (x,y) -> xy, given a tame B containing the inverses), then, under
% « R_0 est un sous-corps de R », a tame Y defined as in (M6) by equations
% f_alpha = e_alpha is rewritten as the common zero set of the f_alpha, and
% with a family (g_beta) embedding X in R^J the products f_alpha g_beta are
% introduced. Page 23: these vanish on Y and separate the points of X \ Y
% (argument written out); a Proposition, under M1 : M7 and M11, that a tame
% Y which is the common zero set of a finite family of tame functions is
% the zero set of a finite family separating the points of X ... X/Y a tame
% space (presumably the quotient; the words between are lost); « Dém : » left empty.
% Page 24: « À prouver » — 1) for a morphism X -> Y of tame spaces, a rare
% tame Y' in Y such that over each connected component U_i of Y \ Y' the
% space X is trivially fibred in the relative tame sense, through a tame
% correspondence R_i in (\bar U_i x Z_i) x X; then the question whether a
% finite diagram of tame spaces over Y is « trivial par morceaux » off a rare
% Y', « Ce serait déjà intéressant » for Z c X c Y x [0,1].
%
% Notation fixed once for the batch: R and R_0 as written, plain italic (no
% blackboard bold: he writes a plain R); his script M for the class of tame
% parts as \mathcal{M}; the axioms as (M8), M6, M1, M7, M11 as he writes them,
% the parentheses only where he has them; the hollow-stroke D opening p. 21
% as \mathbb{D}, with a note; X \smallsetminus Y for his backslash, X/Y where
% he writes a slash (p. 23, the quotient); \pi_0 for the set of connected
% components; « ind-modéré » for the word he writes « ind- » (the first
% stroke reads « rd », « rl »; the reading ind- is suggested by R^I with I
% arbitrary, an ind-object, and is flagged once).
%
% Readings to check first: « ind-modéré » (p. 21, three occurrences); the
% underlined « ensembles » followed by « anneaux » (p. 21) — underline or
% strike?; the interlinear addition of p. 21 (« Dans les applications
% R_0-polynomiales ... transforment des parties modérées en itou »); « la
% restriction » (p. 21 foot); « qui soit ... des M6 comme » (p. 22); the
% words lost in the Proposition of p. 23 (« qui est ... des zéros communs »,
% and the word before X/Y); the last word of p. 24, before « au dessus de Y ».
% Apparatus: 10 \ill{}, 11 \uncertain{}.

\documentclass[11pt,a4paper]{article}
% Le préambule commun à toutes les transcriptions du fonds Grothendieck.
%
% Il tient en une idée : **l'incertitude se compose, elle ne se résout pas.**
% Une transcription qui lisse un mot douteux a détruit la seule chose pour
% laquelle on la faisait — un lecteur doit pouvoir savoir, à chaque endroit,
% ce qui était sur le feuillet et ce qui a été ajouté. Les six macros qui
% suivent sont donc l'appareil critique, et non de la décoration.
%
% Les mêmes commandes sont reconnues par `scripts/render.mjs`, qui produit la
% vue de lecture du site. Une macro ajoutée ici doit l'être aussi là-bas,
% faute de quoi le rendu échouera bruyamment — ce qui est voulu : un rendu
% silencieusement faux est pire qu'un rendu absent.

\ProvidesPackage{grothendieck}[2026/08/08 Transcriptions du fonds Grothendieck]

\usepackage{amsmath,amssymb,amsthm}
\usepackage{mathrsfs}
\usepackage{stmaryrd}
% Les diagrammes commutatifs sont partout dans ce fonds : c'est la langue
% même de la géométrie algébrique de Grothendieck, et une transcription qui
% les rendrait en prose ne transcrirait rien.
\usepackage{tikz-cd}
% Français par défaut — c'est la langue des feuillets — mais l'anglais est
% chargé aussi : la traduction et le résumé sont des documents anglais, et la
% typographie française (espaces avant les deux-points) y serait une faute.
% Ces documents basculent par \selectlanguage{english} en tête de corps.
\usepackage[english,french]{babel}
\usepackage{fontspec}
\usepackage{geometry}
\geometry{a4paper,margin=25mm,marginparwidth=28mm}
\usepackage{marginnote}
\usepackage{xcolor}
% ulem, not soul. `soul`'s \st and \ul are built on a character-by-character
% scan that XeLaTeX's font machinery breaks: the document fails with an
% undefined control sequence rather than degrading. ulem does the same two jobs
% and is Unicode-engine safe. `normalem` stops it hijacking \emph, which the
% transcriptions use for Grothendieck's own emphasis.
\usepackage[normalem]{ulem}
\usepackage{hyperref}
\hypersetup{colorlinks=true,linkcolor=black,urlcolor=black,pdfborder={0 0 0}}

% --- Métadonnées du lot -----------------------------------------------------
% Reprises telles quelles par le script de rendu, qui les affiche en tête de
% la vue de lecture. Elles fixent ce que le fichier prétend couvrir : sans
% elles, une transcription flotte sans point d'attache dans un fonds de seize
% mille feuillets.

\newcommand{\folder}[1]{\gdef\grothFolder{#1}}
\newcommand{\batch}[1]{\gdef\grothBatch{#1}}
\newcommand{\leaves}[2]{\gdef\grothFirst{#1}\gdef\grothLast{#2}}
\newcommand{\foldertitle}[1]{\gdef\grothFoldertitle{#1}}
\gdef\grothFolder{?}\gdef\grothBatch{?}\gdef\grothFirst{?}\gdef\grothLast{?}\gdef\grothFoldertitle{}

% --- Le résumé --------------------------------------------------------------
%
% Il ouvre la lecture modernisée, sous le titre « Résumé », et il est composé
% en petit. Ce n'est pas un ornement : c'est le seul passage du document qui
% ne soit pas des mathématiques, et un lecteur doit pouvoir le reconnaître
% avant de l'avoir lu — puis le sauter s'il connaît déjà le sujet, ou s'y
% arrêter s'il ne le connaît pas. Le filet qui le suit marque la reprise du
% corps.

\newenvironment{resume}
  {\small\setlength{\parskip}{0.45em}}
  {\par\medskip\hrule\medskip}

% --- Mots-clés ---------------------------------------------------------------
%
% En anglais, en clôture du résumé de la lecture modernisée : le vocabulaire
% moderne sous lequel chercher ce que les feuillets construisent. Le manifeste
% du site (`npm run manifest`) les extrait pour étiqueter la cote entière —
% cette ligne est la source unique des tags, il n'y a pas de fichier à côté.
\newcommand{\keywords}[1]{\par\smallskip\noindent
  {\footnotesize\textsc{Keywords} --- \textit{#1}}\par}

% --- L'appareil critique ----------------------------------------------------

% Le numéro de feuillet, en marge. C'est l'ancre qui permet de régler un
% désaccord sur une formule en regardant *un* feuillet plutôt qu'une chemise
% de six cents. Le site s'en sert aussi pour tourner les pages du fac-similé
% quand on fait défiler la transcription.
\newcommand{\leaf}[1]{%
  \marginnote{\footnotesize\color{gray!70}\textsf{#1}}%
  \label{leaf:#1}\ignorespaces}

% Illisible : on ne devine pas. Un mot inventé qui se lit comme les autres est
% le pire résultat possible.
\newcommand{\ill}{\textcolor{red!60!black}{[\,\ldots\,]}}

% Lecture douteuse : le mot est donné, et signalé comme incertain.
\newcommand{\uncertain}[1]{\uline{#1}}

% Ajout de l'éditeur — une lettre manquante, un mot sous-entendu.
\newcommand{\add}[1]{\textcolor{blue!50!black}{[#1]}}

% Biffé par Grothendieck lui-même. Ce qu'il a rayé fait partie du document :
% une rature dit souvent où la pensée a changé de direction.
\newcommand{\struck}[1]{\sout{#1}}

% Note du transcripteur, sur l'état matériel du feuillet.
\newcommand{\note}[1]{\par\smallskip{\small\color{gray!60!black}\textit{[#1]}}\par\smallskip}

% Note marginale de Grothendieck, distincte du corps du texte.
\newcommand{\marginal}[1]{\par\smallskip\noindent\hspace{1em}%
  {\small\color{gray!60!black}#1}\par\smallskip}

% --- Titre ------------------------------------------------------------------

\AtBeginDocument{%
  \begin{center}
    {\large\bfseries Fonds Alexandre Grothendieck --- cote n\textsuperscript{o}~\grothFolder}\\[2pt]
    {\itshape\grothFoldertitle}\\[4pt]
    {\small feuillets \grothFirst--\grothLast\ (lot \grothBatch)}\\[2pt]
    {\footnotesize\color{gray!60!black}Transcription établie d'après le fac-similé numérisé
     par l'Université de Montpellier.\\
     Les lectures incertaines sont soulignées, les illisibles notées [\,\ldots\,],
     les ajouts entre crochets.}
  \end{center}
  \vspace{1em}}

\endinput


\folder{120}
\batch{2}
\pages{21}{24}
\dating{[à partir de 1973-à partir de 1978]}
\watermark{Édition de démonstration}
\foldertitle{Topologie modérée : notes manuscrites (s.d.)}

\begin{document}

\page{21}

$\mathbb{D}$ $=$ pour $X$ un espace \uncertain{ind}-modéré, p.\,ex. (si $R$ est réunion \struck{des} de ses parties modérées (M8),
i.e. \struck{sont} les $R^I$ des espaces ind-modérés) pour les
fonctions ind-modérées sur les $R^I$, qui contiennent en tout cas les
\emph{ensembles} \uncertain{anneaux} de polynômes
$R_0[(X_i)_{i \in I}]$ en les fonctions coordonnées.
\note{$\mathbb{D}$ : une lettre D à double trait, entourée d'une boucle, en
tête du feuillet, suivie du signe $=$ ; sa fonction n'est pas claire sur ces
quatre pages. --- « ind- » : le préfixe s'écrit comme « rd- » ou « rl- » ; la
lecture \emph{ind-} est suggérée par les $R^I$ pour $I$ quelconque qui
suivent, et vaut pour les trois occurrences de la page. --- « les $R^I$ » est
écrit au-dessus de la ligne, en remplacement du mot biffé « sont ». ---
« ensembles » : le trait passe sous le mot, soulignement ou biffure ; le mot
suivant le remplace peut-être}
\marginal{Dans les applications $R_0$-polynomiales \struck{\ill{}}
$R^I \to X$ modérées, et \uncertain{transforment} des parties
\uncertain{modérées} en \uncertain{itou}.}\note{ajout interligne, serré au
droit de « fonctions coordonnées » et au-dessus de l'énoncé suivant ; le mot
biffé en fin de deuxième ligne commence comme « transf », un mot illisible
écrit au-dessus}

\textbf{Corollaire} Soit $f$ une fonction modérée sur l'espace modéré $X$.
Alors $f$ est inv.\ dans l'anneau des fonctions modérées ssi $f$ ne s'annule
pas dans $X$.

En considérant $f(X) \subset R^{*}$, cela revient au

\textbf{Corollaire} Soit $A \subset R^{*}$, $A \in \mathcal{M}$, alors la
\uncertain{restriction} de $x \mapsto x^{-1}$ : $A$ est modérée, i.e.
l'ens.\ des $(x, x^{-1}) \in R^2$, $x \in A$ est $\in \mathcal{M}$.

\page{22}
\note{en haut à droite, « 11 », vraisemblablement sa propre numérotation du
feuillet}

En effet, c'est aussi l'ens.\ des
\[
  \lbrace (x,y) \in R^2 \mid xy = 1,\ x \in A \rbrace .
\]
Si on suppose qu'il existe \struck{$B \subset R^{*}$}, $B \in \mathcal{M}$,
contenant l'ens.\ \struck{\ill{}} $\lbrace x^{-1} \mid x \in A \rbrace$, on
gagne, car l'ens.\ est la partie modérée de $A \times B$ image inverse de
$1 \in A.B$ par l'application modérée $(x,y) \mapsto xy$.
\marginal{$A.B \subset R$, $\in \mathcal{M}$}\note{écrit en colonne sous
« $A.B$ », au bord droit ; lecture incertaine : on lit un signe $=$ ou $''$,
puis $\subset$, $\cap$, et $\mathcal{M}$}

\struck{\ill{}} Cor\note{« Cor » souligné, écrit au-dessus d'un mot
raturé} \quad $R_0$ est un sous-corps de $R$\note{toute la ligne est
soulignée d'un long trait}

Si $X$ modéré, $Y \subset X$ une partie modérée qui soit \ill{}
\uncertain{dans un} des M6 comme
\[
  Y = \bigcap_{\alpha \in I} f_\alpha^{-1}(\lbrace e_\alpha \rbrace),
  \qquad e_\alpha \in R ,
\]
alors quitte à remplacer les $f_\alpha$ par $f_\alpha - e_\alpha$, OPS les
$e_\alpha = 0$, i.e.\ $Y$ est l'ens.\ des \emph{zéros} communs des
$f_\alpha$. Soient $(g_\beta)_{\beta \in J}$ \struck{les} des fonctions
modérées\note{« modérées » est écrit au-dessus, rattaché par un
trait} \uncertain{réalisant} un plongement $X \subset R^J$. Considérons les
$f_\alpha$, $f_\alpha g_\beta$ --- elles sont toutes nulles sur $Y$, et

\page{23}

\struck{si $x \in X \smallsetminus Y$, $\exists\, \alpha$ tel que
$f_\alpha(x) \neq 0$, et} définissent $Y$ comme l'ens.\ des zéros communs.
\struck{de $X$} : en particulier, de plus elles séparent les pts de
$X \smallsetminus Y$, car si $x, y$ sont deux pts distincts,
\struck{$\exists\, \beta$ tel que $f_\alpha(x) \neq 0$
$g_\beta(x) \neq g_\beta(y)$ d'où}

Si les $f_\alpha$ \struck{\ill{}} suffisent pas à les séparer, montrons que
les $f_\alpha g_\beta$ les séparent : \uncertain{en effet}, $\exists\, \beta$
tel que $g_\beta(x) \neq g_\beta(y)$ \struck{d'où} et $\alpha$ tel que
$\lambda = f_\alpha(x) \neq 0$, alors $f_\alpha(y)$ ($= f_\alpha(x) =
\lambda$ par hypothèse) $\neq 0$, d'où
\[
  \underbrace{f_\alpha}_{\lambda} g_\beta(x) = \lambda\, g_\beta(x)
  \neq f_\alpha g_\beta(y) = \lambda\, g_\beta(y) .
\]

Dém :\note{rien ne suit}

\textbf{Proposition} Si on a M1 : M7 et M11, alors pour toute partie
\struck{\uncertain{fermée}} modérée $Y$ d'un esp.\ modéré $X$,
qui est \ill{} des zéros communs d'une famille finie de fonctions modérées
sur $X$, on peut trouver une telle famille qui sépare les pts de $X$
\ill{} $X/Y$ est un espace modéré).
\note{« M1 : M7 » : les deux-points sont tels qu'écrits, peut-être
« M1 à M7 ». --- « modérée » est écrit au-dessus du mot raturé, dont la
lecture « fermée » est elle-même douteuse. --- La parenthèse ouvrante manque ;
le mot illisible avant $X/Y$, en fin de ligne, commence par une capitale}

\page{24}

À prouver

1) Si $X \to Y$ morphisme d'espaces modérés, $\exists$ partie modérée $Y'$
\emph{rare} dans $Y$ et
\[
  (Z_i)_{i \in \pi_0(Y \smallsetminus Y')}, \qquad
  \underbrace{\pi_0(Y \smallsetminus Y')}_{I} = (U_i)_{i \in I},
  \qquad Z_i \ \text{modérés}
\]
\struck{et des} (de sorte que les $\overline{U}_i$ sont des parties modérées
de $Y$) et enfin $\forall\, i \in I = \pi_0(Y \smallsetminus Y')$, on ait
une correspondance modérée
\[
  R_i \subset (\overline{U}_i \times Z_i) \times X
\]
induisant une bijection entre $U_i \times Z_i$ et $X | U_i$, i.e.\
$X | \struck{\ill{}}\, U_i$ est fibré \struck{\ill{}} trivial au sens
modéré relatif ($U_i$ étant un espace modéré relatif \ldots)

Énoncés analogues pour un diagramme fini d'espaces modérés au dessus de
$Y$ ?? Ce diagramme est-il « trivial par morceaux » au dessus d'un
$Y \smallsetminus Y'$ convenable, avec $Y'$ modéré rare dans $Y$ ??? Ce
serait déjà intéressant dans le cas particulier \uncertain{d'un} diagramme
de la forme $Z \subset X \subset Y \times [0,1]$ \ill{} au dessus
de $Y$\ldots
\note{le crochet ouvrant de $[0,1]$ est tracé par-dessus un premier signe,
peut-être un $I$ ; le mot illisible qui suit commence par une capitale
surchargée}

\end{document}
